\section{Disciple, failure, and Easter witness}\label{sec:peter-disciple}

The canonical narratives agree that Peter stood near the center of Jesus'
public ministry. In the Synoptic lists he is named first; Matthew even adds
``first'' before Simon called Peter. With James and John he forms an inner
group present at the raising of Jairus's daughter, the Transfiguration, and
Gethsemane. He speaks when others remain silent, asks questions that open
teaching, and regularly misunderstands the answer. This concentration does
not make him the hero of every scene. It makes his formation unusually
visible.

\subsection{Learning to follow}

Mark's Peter sees a household healing at the beginning of the ministry and
then watches Capernaum gather at the door. He searches for Jesus when prayer
has taken him away from the crowd. He asks about a parable, speaks during the
Transfiguration, notices the withered fig tree, and joins the questions about
the Temple's future. Matthew develops additional scenes: Peter attempts to
walk on the water, asks how often a brother must be forgiven, and receives a
lesson on royal taxation. Luke gives him the miraculous catch and places him
among those asking whether a parable is for the disciples or for all. John
shows Andrew bringing him to Jesus, Peter resisting and then demanding the
washing of his whole body, and later asking about the beloved disciple's
future.

These memories have a coherent moral texture without forming a modern
personality file. Peter is receptive and resistant, courageous and afraid,
generous and self-protective. His quick speech can confess a truth he has not
yet learned to inhabit. Hagiography has sometimes reduced this to an amiable
temperament---the impulsive fisherman---but the evangelists are making a
deeper claim. Apostolic witness is learned under correction.

\subsection{The confession and the rock}

At Caesarea Philippi the Synoptic tradition turns on Peter's confession.
Mark has him say, ``You are the Messiah''; Luke, ``the Messiah of God'';
Matthew, ``the Messiah, the Son of the living God.'' All three immediately
move toward Jesus' prediction of rejection and death. Peter rebukes Jesus,
and Jesus rebukes Peter as Satan for thinking in human rather than divine
terms. The confession is therefore neither cancelled by the correction nor
permission to detach messiahship from the cross.

Matthew alone preserves the extended blessing: flesh and blood have not
revealed the confession; Simon is Peter, and upon this rock Christ will build
his Church; the gates of Hades will not prevail; Peter receives the keys of
the kingdom and authority to bind and loose (Matt 16:17--19). Catholic
tradition reads the commission as personal and enduring: Peter is constituted
the visible rock and steward within a Church that belongs to Christ. The
wordplay and the singular address resist an interpretation in which Peter
simply disappears. At the same time, Matthew later speaks of binding and
loosing to the disciples (18:18), and the New Testament also applies
foundation imagery to the apostles together and supremely to Christ. Petrine
service is relational, not solitary possession.

Luke gives a complementary commission inside the Passion: Satan has demanded
to sift all of the disciples, Jesus has prayed specifically for Simon that his
faith not fail, and, once turned back, he must strengthen his brothers (Luke
22:31--32). John 21 entrusts Jesus' lambs and sheep to him three times. These
texts do not use identical images or arise in identical narrative settings.
Together they explain why Catholic doctrine speaks both of Peter's singular
office and of a college of apostles whose head he is.

\subsection{The Passion: sword, flight, and denial}

At the meal Peter promises fidelity. In Gethsemane he cannot keep watch. All
four Gospels remember an armed blow against the high priest's slave; John
alone identifies the swordsman as Peter and the slave as Malchus. Jesus stops
the violence, though the sayings differ: the cup given by the Father, the
sheathing of the sword, healing, and the warning that those who take the sword
perish by it. Peter's defense does not become apostolic policy.

He then follows at a distance and denies association with Jesus three times.
The precise speakers and wording vary, but the triple failure, the cockcrow,
and bitter recognition are shared. Luke adds Jesus' look; John places the
denials beside a charcoal fire. The tradition did not hide its foremost
witness's collapse. Nor did it treat the collapse as his final identity. The
same Peter who said he would die rather than deny must receive his ministry as
forgiven vocation, not as the reward for having proved stronger than the
others.

\subsection{First among named witnesses}

The earliest datable Easter reference to Peter is not a Gospel narrative but
the formula Paul says he received: Christ died, was buried, was raised, ``and
that he appeared to Cephas, then to the Twelve'' (1 Cor 15:3--5). The formula
does not describe the appearance, but its placement gives Cephas a primary
role in the public transmission of Resurrection faith. Luke 24:34 similarly
reports that the Lord has risen and appeared to Simon. Mark's empty-tomb
messenger singles out ``his disciples and Peter.'' Matthew emphasizes the
commission to the eleven in Galilee. John 20 places Peter and the beloved
disciple at the tomb; John 21 narrates a lakeside appearance culminating in
the threefold charge to feed the flock and a foreshadowing of the kind of
death by which Peter would glorify God.

The canonical ensemble is not a contest over who loved Jesus most. Mary
Magdalene's indispensable first proclamation, the beloved disciple's insight,
the apostolic group, and Peter's public commission coexist. Peter's authority
is paschal: the denied Lord calls the denier to love, service, and ultimately
the cross.

\begin{studybox}{Apostolic sanctity}
Peter's failures are not late hostile inventions that the Church reluctantly
tolerated. They are woven into all four canonical portraits. The Church
venerates the restored witness, not an imaginary man who never needed mercy.
\end{studybox}
